\section{Mecanismo de Encapsulamiento de Claves ML-KEM (FIPS 203)}
% =============================================================================

En agosto de 2024, el NIST (2024a) publicó el estándar oficial \textbf{FIPS 203}, el cual define el algoritmo de encapsulamiento de claves basado en retículos modulares \textbf{ML-KEM}, originado en el proyecto CRYSTALS-Kyber. Este estándar constituye el mecanismo principal a nivel mundial para proteger la confidencialidad de la información frente a computadoras cuánticas.

\subsection{Concepto de KEM en Votación Electrónica}

Cuando se estudia criptografía por primera vez, suele pensarse que la clave pública sirve para cifrar directamente archivos o mensajes largos. Sin embargo, en la práctica informática moderna se utiliza una técnica mucho más eficiente denominada \textbf{Mecanismo de Encapsulamiento de Claves} (KEM).

El funcionamiento de un KEM puede explicarse mediante una analogía cotidiana: supóngase que el votante desea enviar su boleta electoral dentro de una caja fuerte muy rápida y ligera (cifrado simétrico AES-256-GCM), para lo cual inventa en ese instante una llave secreta única $K$ de $256$ bits. El problema es cómo hacer llegar esa llave $K$ a la Autoridad de Conteo a través de Internet sin que ningún atacante la copie en el camino.

Para resolverlo, el algoritmo ML-KEM no cifra el voto directamente, sino que ``encapsula'' (envuelve) esa llave simétrica $K$ utilizando la clave pública post-cuántica de la autoridad electoral, produciendo una cápsula cifrada $c$. De este modo, el votante cifra su voto con la llave $K$ mediante AES-256-GCM y adjunta la cápsula $c$. Al final de la elección, solo la Autoridad de Conteo posee la clave privada necesaria para abrir la cápsula $c$, recuperar la llave $K$ y leer el contenido del voto.

\subsection{Fases del Protocolo ML-KEM}

El funcionamiento de ML-KEM se compone de tres pasos secuenciales bien definidos: generación de claves, encapsulamiento y desencapsulamiento. En la Figura \ref{fig:kem_flow} se presenta el diagrama general de interacción entre la autoridad electoral y la terminal del votante.

\begin{figure}[!htbp]
	\centering
	\begin{tikzpicture}[
		pbox/.style={rectangle, draw=black, fill=white, line width=0.8pt, rounded corners=1.5pt, text width=3.7cm, align=center, font=\footnotesize\rmfamily, inner sep=5pt},
		decbox/.style={rectangle, draw=black, fill=white, line width=0.8pt, rounded corners=1.5pt, text width=4.2cm, align=center, font=\footnotesize\rmfamily, inner sep=5pt},
		arr/.style={-{Latex[length=2.2mm, width=1.5mm]}, line width=0.8pt, draw=black}
	]
		% Servidor Electoral y Terminal del Votante con amplia separación horizontal
		\node (aut) [pbox] at (-4.4, 0) {
			\textbf{Servidor Electoral}\\[3pt]
			\scriptsize Genera $(\mathbf{A}, \mathbf{b})$ con\\vector secreto $\mathbf{s}$
		};
		\node (vot) [pbox] at (4.4, 0) {
			\textbf{Terminal del Votante}\\[3pt]
			\scriptsize Encapsula secreto $K$\\en $c = (\mathbf{u}, v)$
		};
		
		% Flechas de intercambio con etiquetas centradas sin superposición
		\draw[arr] ([yshift=4mm]aut.east) -- ([yshift=4mm]vot.west) 
			node[midway, above=2pt, font=\scriptsize\rmfamily] {1. Clave pública $PK = (\mathbf{A}, \mathbf{b})$};
		
		\draw[arr] ([yshift=-4mm]vot.west) -- ([yshift=-4mm]aut.east) 
			node[midway, below=2pt, font=\scriptsize\rmfamily] {2. Boleta cifrada $c = (\mathbf{u}, v)$};
		
		% Desencapsulamiento
		\node (dec) [decbox] at (-4.4, -2.3) {
			\textbf{Desencapsulamiento}\\[3pt]
			\scriptsize Calcula $v - \mathbf{s}^T\mathbf{u}$ y\\redondea para recuperar $K$
		};
		\draw[arr] (aut) -- (dec);
	\end{tikzpicture}
	\caption{Flujo de operaciones del algoritmo ML-KEM.}
	\label{fig:kem_flow}
	\par\vspace{3pt}{\footnotesize\textit{Fuente: Adaptado de NIST (2024a).}}
\end{figure}

Tal como se observa en la Figura \ref{fig:kem_flow}, la autoridad pone a disposición la clave pública $PK = (\mathbf{A}, \mathbf{b})$, el terminal del votante la utiliza para encapsular el secreto $K$ dentro del criptograma $c = (\mathbf{u}, v)$, y posteriormente se realiza el desencapsulamiento restando el término secreto $\mathbf{s}^T\mathbf{u}$. A continuación se explica de forma sencilla qué ocurre en cada una de estas tres fases:

\begin{enumerate}
	\item \textbf{Generación de Claves:}
	Antes de iniciar la votación, la autoridad genera una matriz pública de polinomios $\mathbf{A} \in R_q^{k \times k}$, un vector de polinomios secretos con coeficientes pequeños $\mathbf{s}$ y un vector de ruido pequeño $\mathbf{e}$. Con ellos calcula el vector público con ruido:
	\begin{equation}
		\mathbf{b} = \mathbf{A}\mathbf{s} + \mathbf{e} \pmod q
	\end{equation}
	La clave pública publicada es el par $PK = (\mathbf{A}, \mathbf{b})$, mientras que la clave privada guardada en custodia es $SK = \mathbf{s}$.

	\item \textbf{Encapsulamiento:}
	Cuando un ciudadano va a emitir su voto, el sistema toma la clave pública $PK = (\mathbf{A}, \mathbf{b})$ y genera un vector aleatorio temporal $\mathbf{r}$ junto con nuevos ruidos pequeños $\mathbf{e}_1$ y $e_2$. A partir de una semilla aleatoria $m$ de $256$ bits (de la cual se deriva la llave simétrica $K$), calcula el texto cifrado compuesto por dos partes $c = (\mathbf{u}, v)$:
	\begin{align}
		\mathbf{u} &= \mathbf{A}^T \mathbf{r} + \mathbf{e}_1 \pmod q \\
		v &= \mathbf{b}^T \mathbf{r} + e_2 + \text{Codificar}(m) \pmod q
	\end{align}
	Aquí la función $\text{Codificar}(m)$ realiza una tarea muy ingeniosa para proteger cada bit de información frente al ruido: como el módulo es $q = 3329$, si el bit de $m$ es $0$ lo coloca en el valor $0$, pero si el bit es $1$ lo coloca a mitad del reloj modular, en $\lfloor 3329 / 2 \rceil = 1665$. Al estar tan separados el $0$ y el $1665$, un pequeño ruido de $\pm 10$ o $\pm 20$ no logrará confundir un uno con un cero.

	\item \textbf{Desencapsulamiento:}
	Durante el escrutinio, la autoridad utiliza su clave privada $\mathbf{s}$ para multiplicar $\mathbf{s}^T \mathbf{u}$ y restarlo de $v$:
	\begin{align}
		v - \mathbf{s}^T \mathbf{u} &= \left( \mathbf{b}^T \mathbf{r} + e_2 + \text{Codificar}(m) \right) - \mathbf{s}^T \left( \mathbf{A}^T \mathbf{r} + \mathbf{e}_1 \right) \nonumber \\
		&= \left( (\mathbf{A}\mathbf{s} + \mathbf{e})^T \mathbf{r} + e_2 + \text{Codificar}(m) \right) - \left( \mathbf{s}^T \mathbf{A}^T \mathbf{r} + \mathbf{s}^T \mathbf{e}_1 \right) \nonumber \\
		&= \text{Codificar}(m) + \underbrace{\left( \mathbf{e}^T \mathbf{r} + e_2 - \mathbf{s}^T \mathbf{e}_1 \right)}_{\mathbf{e}_{\text{residual}}} \pmod q
	\end{align}
\end{enumerate}

Como muestra el resultado final de la ecuación anterior, los términos grandes $\mathbf{s}^T \mathbf{A}^T \mathbf{r}$ se anulan entre sí de manera exacta. Lo único que queda es el mensaje $\text{Codificar}(m)$ sumado a un ruido residual pequeño $\mathbf{e}_{\text{residual}}$. Finalmente, el programa aplica una operación de redondeo: si el número resultante está cerca de $0$, recupera el bit $0$; si está cerca de $1665$ (la mitad de $q$), recupera el bit $1$, reconstruyendo sin errores los $256$ bits de la clave $K$ que permite abrir el voto.

\subsection{Seguridad IND-CCA2 y Transformada Fujisaki-Okamoto}

En ciberseguridad no basta con proteger el sistema frente a un atacante que solo escucha la red; también es necesario defenderse de un atacante activo que modifique ligeramente los textos cifrados enviados al servidor para observar si el sistema devuelve mensajes de error y así tratar de adivinar la clave privada. La resistencia contra este tipo de ataques se conoce formalmente como \textbf{indistinguibilidad frente a ataques de texto cifrado escogido adaptativo} o seguridad IND-CCA2.

Para alcanzar la seguridad IND-CCA2, el estándar FIPS 203 incorpora un mecanismo de verificación interna llamado \textbf{Transformada de Fujisaki-Okamoto} (Fujisaki y Okamoto, 2013). Su funcionamiento consiste en que, inmediatamente después de desencapsular y obtener el valor candidato $m'$, el receptor vuelve a ejecutar internamente el paso de encapsulamiento usando ese mismo $m'$ y comprueba si la cápsula generada coincide bit por bit con la cápsula $c$ que llegó por la red.

Si un atacante alteró aunque sea un solo bit del criptograma en el camino, ambas cápsulas no coincidirán. En ese caso, ML-KEM rechaza silenciosamente el paquete manipulado sin dar pistas del fallo, garantizando la protección total de la clave privada electoral.

% =============================================================================
